\documentclass{beamer}
\usepackage[utf8]{inputenc}

%\usetheme{Copenhagen}
\usetheme{Antibes}
\hypersetup{pdfstartview={Fit}}

\title{Joint 3D object recognition and tracking}
\author[Orizio, Rizzini, Zucchelli]{Orizio Riccardo,\\Rizzini Mattia,\\Zucchelli Maurizio}
\date{26 luglio 2013}
\institute[UniBS]{Universit\`{a} degli Studi di Brescia}
\logo{\includegraphics[width=15mm]{logo.png}}

\begin{document}
	\begin{frame}
		\maketitle
	\end{frame}

	\section{Introduction}
	
	\begin{frame}
		\frametitle{\insertsection}
		The project goal is to realize a joint system of recognition and tracking
		of some features of an object which is identified in a video.
		These features are called \textsc{Label}.\\
		The sofware is aimed for execution on low performance devices such as the
		iPad, therefore the tracking is necessary due to the high computational
		cost of a pure recognition.\\
	\end{frame}

	\begin{frame}
		\frametitle{\insertsection\ - Overall structure}
		The software is designed for concurrent execution.\\
		There is an object of class \textsc{Manager} which, given every frame
		captured by a webcam or in a video, elaborates it in some way and retuns
		an object of class \textsc{Object} which contains the actual \textsc{Label}
		positions.\\
	\end{frame}

	\AtBeginSection[]
	{
		\begin{frame}
			\frametitle{Outline}
			\tableofcontents[currentsection]
		\end{frame}
	}
	\section{Manager}

	\begin{frame}
		\frametitle{\insertsection}
		The \textsc{Manager} handles two objects of class \textsc{Recognizer} and
		\textsc{Tracker}. The first object runs a \emph{thread}.\\
		Every $N$ frames elaborated, the \textsc{Manager} asks the
		\textsc{Recognizer} to perform a full recognition. In the meantime the
		\textsc{Tracker} keeps tracking the \textsc{Label}s starting from the last
		\textsc{Object} given by the \textsc{Recognizer}.\\
		When the \textsc{Recognizer} has finished it returns an \textsc{Object}
		which contains the positions of the \textsc{Label}s in the moment in which
		the recognition started. These positions are \emph{actualized} by the
		\textsc{Tracker}.\\
	\end{frame}

	\section{Recognition}

	\begin{frame}
		\frametitle{\insertsection}
		To perform recognition, \emph{SURF} descriptors are used in place of
		\emph{SIFT} ones because they are proved to be lighter in application.\\
		The recognition bases itself on a \textsc{Database}.\\
		This \textsc{Database} support serialization of the structures used for
		recognition. At the program start, a \textsc{Database} name is provided:
		if this name corresponds to a serialized \textsc{Database}, the
		\textsc{Database} is loaded; otherwise it is trained.\\
			% Maybe eccessive `textsc` for the word database? Lol.
	\end{frame}

	\subsection{Database}

	\begin{frame}
		\frametitle{\insertsection\ - \insertsubsection}
		The \textsc{Database} training is performed through the analysis of a group
		of sample images depicting the object from various points of view.\\
		To every image is also associated a text file which contains the absolute
		positions of the various labels in the referenced sample.\\
		Every sample image is loaded, from every image \emph{SURF} features are
		extracted, descriptors computed and label positions loaded. Everything
		is then saved in three structures and serialized.\\
	\end{frame}

	\subsection{Matching}

	\begin{frame}
		\frametitle{\insertsection\ - \insertsubsection}
		At the end of the training or load phase, a \emph{FLANN} based matcher is
		trained based on all the samples descriptors.\\
		When the \textsc{Manager} asks the \textsc{Recognizer} to perform a match
		on a frame, his \emph{SURF} features are extracted and descriptors
		computed. The descriptors are then used to find a match in the
		\emph{FLANN} index.\\
		For every descriptor, the two best matches are given. Then the matches are
		filtered to exclude the ones in which the distance of the best match are
		too close to the distance of the second best match.\\
	\end{frame}

	\begin{frame}
		\frametitle{\insertsection\ - \insertsubsection}
		These good matches are analyzed to find out what is the sample who's got
		more of them. That sample is the only considered and his matching keypoints
		are extracted from the \textsc{Database}.\\
		Then an homography matrix is calculated between the sample keypoints and the
		ones in the frame. The matrix is used to map the positions of the
		\textsc{Label}s in the frame. These \textsc{Label}s are added to an
		\textsc{Object} which is returned to the \textsc{Manager}.\\
	\end{frame}

	\section{Tracking}

	\begin{frame}
		\frametitle{\insertsection}
		The tracking of the object is divided in two principal parts:
		\begin{enumerate}
			\item The effective tracking of the \textsc{Object} between two
				successive frames.
			\item The \emph{actualization} of the \textsc{Object} received by the
				\textsc{Recognizer}.
		\end{enumerate}
	\end{frame}

	\begin{frame}
		\frametitle{\insertsection}
		To track an object, \emph{GFTT} (Good Features To Track) are used:
		these features are extremely fast to calculate and provide optimal results
		in tracking application.\\
		The algorithm works on the whole frame: the features are associated to the
		new frame using Lucas-Kanade's method for \emph{optical flow} estimation,
		then, for every \textsc{Label}, the nearest features to it are calculated
		using a \emph{FLANN} based matcher and their movement between the frames is
		mediated and added to the \textsc{Label}'s position.
	\end{frame}

	\subsection{Actualization}

	\begin{frame}
		\frametitle{\insertsection\ - \insertsubsection}
		The \emph{actualization} process updates the \textsc{Object} received from
		the \textsc{Recognizer} to its position in the current frame. This is
		necessary because of the long time required by it.\\
		To allow an easy and fast \emph{actualization}, the \textsc{Tracker}
		calculates and saves the features of the frame used for the recognition
		and starts using them internally for the tracking of the other frames.\\
		This way, when the \textsc{Object} is received, the only operation to be
		done is the tracking from the saved features to the current features.\\
	\end{frame}
	
	\subsection{Optimization}

	\begin{frame}
		\frametitle{\insertsection\ - \insertsubsection}
		The heaviest part of the algorithm is the calculation of the \emph{optical
		flow}. To reduce the time required, the frame is decimated of a factor
		$2$; since the complexity of the algorithm is higher than linear, this
		reduces the computational cost to less than $25\%$ the original time.\\
		Also, to furtherly simplify the complexity, the whole tracker has been
		adapted to work with decimated frames. This can be done without damaging
		the results because both the features and the tracking algorithm are robust
		to scaling.\\
	\end{frame}

	\section{Conclusions}

	\begin{frame}
		\frametitle{\insertsection}
		It works, so shut up!
	\end{frame}
\end{document}
